\DSource{0133}{44}
\hypertarget{AB01-PDF0133-P01}{}
\DLine{AB01-PDF0133-body-L001}{}{}{partem quam Sol medio suo itinere ab initio Arietis ad initium Librae percurrit, scilicet, ut}
\DLine{AB01-PDF0133-body-L002}{}{}{supra vidimus, $183^\circ56'12''$. Arcus KLRM, semicircumferentia excentrici, $180^\circ$ continet; quis-}
\DLine{AB01-PDF0133-body-L003}{}{}{que igitur arcuum PK, YM est dimidia pars eorum $3^\circ56'12''$, quibus Sol medio suo itinere}
\DLine{AB01-PDF0133-body-L004}{}{}{$180^\circ$ superat, idest $1^\circ58'6''$. --- Manifestum etiam est arcum PKLR (1) excentrici id esse quod}
\DLine{AB01-PDF0133-body-L005}{5}{}{Sol, medio suo itinere, ab initio Arietis ad initium Cancri percurrit; igitur $92^\circ14'10''$ erit;}
\DLine{AB01-PDF0133-body-L006}{}{}{et quia arcus PKL, ob ea quae nuper diximus, $91^\circ58'6''$ complectitur, arcus LR $0^\circ16'4''$}
\DLine{AB01-PDF0133-body-L007}{}{}{continebit. Item manifestum est perpendicularem PQ sinum esse arcus PK, et perpendicula-}
\DLine{AB01-PDF0133-body-L008}{}{}{rem RG sinum arcus LR; igitur PQ fere $2^p3'39''$, et RG fere $0^p16'45''$ erit. Quoniam}
\DLine{AB01-PDF0133-body-L009}{}{}{linea KM parallela est lineae AC, erit linea ES aequalis PQ; et similiter, quia LN parallela}
\DLine{AB01-PDF0133-body-L010}{10}{}{est lineae BD, erit FS aequalis RG. Est itaque etiam latus EF trianguli rectanguli ESF no-}
\DLine{AB01-PDF0133-body-L011}{}{}{tum; quadratum lateris ES est fere $4^p14'48''$, quadratum lateris FS fere $0^p4'41''$, summa}
\DLine{AB01-PDF0133-body-L012}{}{}{igitur $4^p19'29''$ (2) efficit quadratum lateris EF, cuius radix $2^p4'\,\qfrac{3}{4}$ est quantitas lineae EF}
\DLine{AB01-PDF0133-body-L013}{}{}{ambo centra coniungentis. Ea ratione autem, qua quadrans circuli triangulum rectangulum}
\DLine{AB01-PDF0133-body-L014}{}{}{ESF circumdantis $90^\circ$ complectitur, et semidiametrus $60^p$, erit arcus EF circiter $1^\circ59'$ (3); et}
\DLine{AB01-PDF0133-body-L015}{15}{}{haec est maxima anomalia motus Solis per has observationes manifesta ({\itshape b}).}
\hypertarget{AB01-PDF0133-P02}{}
\DLine{AB01-PDF0133-body-L016}{}{}{\Indent Nunc in quantitatem arcus BH eclipticae inquiramus; qua cognita, etiam arcus HA reli-}
\DLine{AB01-PDF0133-body-L017}{}{p. 67.}{quus notus erit. Punctum T est locus apogei in circulo Solis excentrico; nam, si lineam EF,}
\DLine{AB01-PDF0133-body-L018}{}{}{ambo centra coniungentem, usque ad eclipticam producimus, ipsa circulum KLMN in T, et}
\DLine{AB01-PDF0133-body-L019}{}{}{circulum eclipticae in H secat. Necesse est ut rationem lineae EF ad semidiametrum EH co-}
\DLine{AB01-PDF0133-body-L020}{20}{}{gnoscamus, nec non quantitatem arcus BH eclipticae. Iam vidimus lineam EF $2^p4'\,\qfrac{3}{4}$ com-}
\DLine{AB01-PDF0133-body-L021}{}{}{plecti ea ratione qua semidiametrus $60^p$ continet; eadem ratione igitur linea EH, quae, ut EB,}
\DLine{AB01-PDF0133-body-L022}{}{}{semidiametrus est, lineam EF 28 vicibus et $\qfrac{5}{6}$ fere continebit. Quoniam linea FS, ut dixi-}
\DLine{AB01-PDF0133-body-L023}{}{}{mus, [$0^p16'45''$] est, si linea EF $60^p$ ponatur, erit iuxta hanc rationem linea FS circiter $8^p4'$;}
\par\noindent\begin{minipage}[t]{0.46\FullSourceWidth}\vspace{0pt}\centering\hypertarget{AB01-PDF0133-F01}{}\includegraphics[width=\linewidth]{AB01-PDF0133-F01.png}\end{minipage}\hfill\begin{minipage}[t]{0.5\FullSourceWidth}\vspace{0pt}\setlength{\GutterOffset}{0.5\FullSourceWidth}
\DLine{AB01-PDF0133-body-L024}{}{}{hoc enim invenitur si $28\,\qfrac{5}{6}$ vicibus [$0^p16'45''$]}
\DLine{AB01-PDF0133-body-L025}{25}{}{multiplicantur. --- Aut, si mavis, lineam FS in}
\DLine{AB01-PDF0133-body-L026}{}{}{semidiametrum EH multiplica; invenies $16^p45'$,}
\DLine{AB01-PDF0133-body-L027}{}{}{quae, per lineam EF scilicet per $2^p4'\,\qfrac{3}{4}$ divisa,}
\DLine{AB01-PDF0133-body-L028}{}{}{dabunt $8^p4'$, quod est [sinus] anguli BEH; ar-}
\DLine{AB01-PDF0133-body-L029}{}{}{cus igitur BH $7^\circ43'$ circiter continet. Itaque}
\DLine{AB01-PDF0133-body-L030}{30}{}{manifestum est punctum T apogei in excentrico,}
\DLine{AB01-PDF0133-body-L031}{}{}{a puncto solstitii aestivi versus partem anterio-}
\DLine{AB01-PDF0133-body-L032}{}{}{rem signorum $7^\circ43'$ distare, scilicet $82^\circ17'$ ab}
\DLine{AB01-PDF0133-body-L033}{}{}{initio Arietis esse. Hoc invenire volebamus. Ob-}
\DLine{AB01-PDF0133-body-L034}{}{}{servatio iuxta quam hunc calculum fecimus,}
\DLine{AB01-PDF0133-body-L035}{35}{}{anno 1194 aerae Dhū ’l-qarnayn (4) fuit, quo}
\DLine{AB01-PDF0133-body-L036}{}{}{iter Solis ab initio Arietis ad initium Cancri et}
\DLine{AB01-PDF0133-body-L037}{}{}{ad initium Librae observavimus ({\itshape c}).}
\hypertarget{AB01-PDF0133-P03}{}
\DLine{AB01-PDF0133-body-L038}{}{}{\Indent Restat nunc nobis ut hanc anomaliam ad}\end{minipage}\par
\DLine{AB01-PDF0133-body-L039}{}{}{gradus eclipticae distribuamus, et portiones eius quae ad singulos gradus pertinent in tabulis}
\par\vspace{6pt}\hrule height.25pt\vspace{5pt}
\noindent\begin{minipage}[t]{.48\textwidth}\vspace{0pt}\fontsize{9}{11.5}\selectfont\setlength{\GutterOffset}{0pt}
\hypertarget{AB01-PDF0133-N01}{}
\DLine{AB01-PDF0133-notes_left-L001}{}{}{\Indent (1) Male codex PKL.}
\hypertarget{AB01-PDF0133-N02}{}
\DLine{AB01-PDF0133-notes_left-L002}{}{}{\Indent (2) Male secundae apud Platonem $59''$.}
\hypertarget{AB01-PDF0133-N03}{}
\DLine{AB01-PDF0133-notes_left-L003}{}{}{\Indent (3) Plato male $1^\circ58'$; \Name{Ibn Yūnus} (apud \Name{Caus-}}
\DLine{AB01-PDF0133-notes_left-L004}{}{}{\Name{sin}, {\itshape Le livre de la grande table Hakémite}, in}
\DLine{AB01-PDF0133-notes_left-L005}{}{}{Notices et extraits des mss. de la Bibl. Nation., t. VII,}
\DLine{AB01-PDF0133-notes_left-L006}{}{}{Paris 1804, p. 155) scribit: « al-Battānī aequationem}
\DLine{AB01-PDF0133-notes_left-L007}{}{}{« totam Solis calculavit, invenitque $1^\circ59'10''$ esse ».}
\end{minipage}
\hfill\begin{minipage}[t]{.48\textwidth}\vspace{0pt}\fontsize{9}{11.5}\selectfont\setlength{\GutterOffset}{0pt}
\DLine{AB01-PDF0133-notes_right-L001}{}{}{\Name{Halley} p. 917, et \Name{Delambre}, {\itshape Astr. moyen âge},}
\DLine{AB01-PDF0133-notes_right-L002}{}{}{p. 36, qui hodiernis tabulis trigonometricis usi sunt,}
\DLine{AB01-PDF0133-notes_right-L003}{}{}{invenere quoque $1^\circ59'10''$. Codicis lectio igitur bona}
\DLine{AB01-PDF0133-notes_right-L004}{}{}{est; in tabulis aequationis Solis (fol. 191,r.) re vera}
\DLine{AB01-PDF0133-notes_right-L005}{}{}{maxima aequatio est $1^\circ59'10''$.}
\hypertarget{AB01-PDF0133-N04}{}
\DLine{AB01-PDF0133-notes_right-L006}{}{}{\Indent (4) Incipit, iuxta usum auctoris nostri, die 1 Se-}
\DLine{AB01-PDF0133-notes_right-L007}{}{}{ptembris 882 post Chr. n.}
\end{minipage}
